\section{Introduction}
%



It is estimated that almost 50\% of the world's data will reside on
remote storage services, e.g., public clouds, by the end of 2025~\cite{public_storage_vol}.
Unfortunately, storage services are often misused to store harmful
content, e.g., child-sexual-abuse-material (CSAM)~\cite{williams2022csam}. In response,
large storage providers like Apple, Google, and Amazon have their own
mechanisms to identify such harmful content~\cite{apple2021csam,
  googleCSAM, amazonCSAM}. However, the task is more challenging
when the content itself is encrypted under keys that are not available 
to the provider~\cite{content_mod_e2ee}. 

In 2021, Apple proposed a mechanism to detect CSAM in encrypted photos
uploaded to iCloud. The system, called the Apple
PSI~\cite{bhowmick2021apple}, uses Apple's proprietary perceptual
hashing technology, NeuralHash, to generate a hash of the
content. The hash is then compared privately against a blocklist curated by Apple using known content. 
If there is a match, then the iCloud server learns the decryption key. However, 
the system was not deployed due to several concerns~\cite{applePSIabandon}. 
First, NeuralHash is known to be fragile, and can
produce false positives~\cite{neuralHashAttack}; unfortunately, the system was not compatible 
with other kinds of perceptual hash algorithms. 
Second, Apple's involvement in creating the
blocklist was considered problematic as it could be misused to
surveil customers~\cite{applePSICriticBBC,
  applePSICriticBU}. 

An alternative to the Apple PSI system was proposed by \citet{kulshrestha2021identifying} that 
can be adapted to different types of perceptual hash algorithms. However, their schemes do not tolerate a 
malicious server, and require an interactive protocol for every upload 
with non-trivial communication and compute costs for the client. 
This can be problematic for image services where often the client devices are smartphones or tablets with limited resources.



In this work, we ask the following question: 

\begin{myquote}
Can we design an encrypted content moderation system where a group of independent auditors create 
the blocklist such that i) clients can authenticate the blocklist before use, ii) the mechanism is compatible with different types of content hashing algorithms, iii) multiple service providers can use the same blocklist without being involved in its curation, and iv) the mechanism has low compute and communication costs?
\end{myquote}


%While we be mitigated, e.g., by
%using a different perceptual hash, and/or by authenticating the
%blocklist by a set of trusted
%\emph{auditors}~\cite{scheffler2023public}, this system was ultimately
%not deployed.

%In this work, we propose mechanisms for \emph{auditable} content
%moderation for encrypted storage services. 

In particular, we envision
a setting where a group of independent parties (auditors), e.g., law enforcement,
and child safety groups, publish a common (obfuscated) blocklist. Such blocklists are
not uncommon, e.g., the NCEMC and FBI have published lists of missing
children and persons~\cite{ncemc}. We 
expect that a threshold number of these independent auditors will behave honestly, and posit that having a common blocklist curated by independent
agencies with different interests, will build more trust in the system
than individual blocklists published by self-interested service
providers. 

We propose a mechanism whereby a group of
auditors can publish a common blocklist in an unforgeable but publicly verifiable
\emph{blocklist authentication token}. The token hides the blocklist elements from the clients, so 
as to ensure that the checks are not trivially circumvented and victim identities are not 
publicly revealed. During uploads, the \emph{blocklist authentication token} is
used by clients to generate content-specific metadata
in the form of a \emph{content verification token}, which is provided to the servers (along with the encrypted data)
and allows checks against the blocklist. The
\emph{content verification token} is hiding in the sense that if the uploaded content is
not in the blocklist, then no information is revealed to the server.
We realize this design to allow for either \emph{exact matching} of content, where an equality test between
content hashes suffices, or \emph{fuzzy matching}, where a metric
distance is computed between content embeddings to estimate
similarity~\cite{yang2025certphash}.



%all parties
%except for the storage servers. The latter is to keep the blocklist
%confidential since otherwise this information can be used to circumvent checks and/or 
%publicly reveal victim identities. The token is constant-sized independent of the number of items
%in the blocklist. T



%We also allow for blocklists that either require
%



%\myparagraph{Exact matching against blocklist}
%
We design the \emph{blocklist authentication token} using a new cryptographic primitive, a \emph{randomizable set-homomorphic
signature}. Set-homomorphic signatures~\cite{johnson2002homomorphic} enable signing a set of items such that given a
signature and the public key, a party can compute a signature
on any subset of the items. We further extend the definition with a \emph{randomizing} function. Informally, given a signature on a
set \blocklist, applying this function to the homomorphically computed
signature on $\blocklist \setminus \{\genericCInputElmt\}$ ensures
that the resulting value hides $\genericCInputElmt$ when
$\genericCInputElmt \notin \blocklist$, even when $\blocklist$ is
known. However, when $\genericCInputElmt \in \blocklist$, the function
produces a signature on $\blocklist \setminus
\{\genericCInputElmt\}$ that can be verified efficiently. 
%

%\footnote{It is also possible to compute
%the signature on the union of two sets, given the signatures on the
%individual sets, with only the public key, but this property is not important for our construction.}

We use the auditors' signature on the
blocklist \blocklist as the \emph{blocklist authentication token}, thus ensuring unforgeability and public verifiability. 
Then, the client derives the \emph{content verification token} by homomorphically computing the signature
on $\blocklist \setminus \{\genericCInputElmt\}$, and applying the
randomizing function to this signature. If $\genericCInputElmt \in
\blocklist$, the server learns $\genericCInputElmt$ as it can verify
the signature, but otherwise the randomizing function hides
$\genericCInputElmt$ from the server. 

We propose a new construction for a
randomizable set homomorphic signature scheme under the strong QR-RSA assumption that produces a constant-sized signature, \emph{independent 
of the blocklist size}. The scheme supports distributed key generation and allows a group of independent auditors with semi-honest majority to sign and update the blocklist in a distributed fashion when the private key is secret-shared among the auditors. We leverage this construction to design a content moderation protocol with optimal client costs. A key feature of our construction is a \emph{fast path} for identifying safe content that eliminates most of the expensive computation (both for the client and the server) required for a full verification, when an (encrypted) uploaded file is not on the blocklist. Thus, our protocol has modest average case costs, since presumably a large portion of the content uploaded to a server will not be blocklisted. 

%\myparagraph{Fuzzy matching against blocklist}
%
We extend the exact-matching mechanism to compute a fuzzy
match when the content hashes are vectors 
in Hamming space\footnote{Our choice of Hamming distance is informed by the fact that several image embeddings use Hamming distance to identify similar images~\cite{singh2025perceptual,neuralHashAttack,yang2025certphash} and existing work on encrypted content moderation also use Hamming distance as the metric~\cite{kulshrestha2021identifying}}. This is achieved by augmenting the protocol with a relatively inexpensive 
coarse-grained check using projections of the vectors to indicates whether a more expensive check using the \emph{content verification token} 
is required. The coarse-grained check eliminates a large fraction of the true negatives and reduces server costs. Crucially, both the checks require communication and compute costs (server + client) that scale sublinearly in the blocklist size. 


\iffalse
 but the overall scheme does not scale to large blocklists due
to high server costs. Our second technical contribution is a
mechanism that enables fuzzy matching when the content embeddings are
in Hamming space. Essentially, the scheme generalizes the Apple PSI
mechanism to fuzzy matching. In short, the server stores a per-client
cuckoo hash table~\cite{fan2014cuckoo}, indexed by random projections
of the blocklist vectors. The values in the table are the blocklist
vectors themselves. During uploads, a client retrieves some elements
from the hash table based on the projections of its input, and
computes the Hamming distance between its own content embedding and
the results returned by the server. Crucially, the hash table is
obfuscated by the client during a preprocessing step such that the
server does not learn anything from the queries (including the Hamming
distance between the blocklist elements and the client's input) unless
there is a match, and cannot tamper with the entries in the
table. While the preprocessing step is expensive, it is performed only
once. The query costs scale sublinearly in the blocklist size.
\fi 
%\mkr{replaced ``compute time'' with latency ... ok?}


\myparagraph{Evaluation}
%
Our exact-matching scheme is highly efficient. The fast path verification takes less than 8ms and a full verification 
takes roughly 430 ms for a blocklist with million elements. Assuming that realistically less than 10\% of the content would require a full check, because uploading CSAM is atypical user behavior, our protocol on average takes 43ms, \textbf{12$\times$ faster} than the PEMC protocol of \citet{kulshrestha2021identifying} \emph{that operates under a more restrictive threat model with a semi-honest server}. Our protocol also has \textbf{50$\times$} lower communication costs. 
While the Apple PSI system does not report performance, 
we estimate it has similar throughputs but with higher client storage costs and produces false negatives. 

%It takes less
%than 8 milliseconds for the client to compute the \emph{content verification token}, and requires less
%than 1KB of communication (excluding the actual encrypted content),
%\emph{independent of the blocklist size}. For a blocklist of a million elements, the server requires on an average 300 ms to check content. The scheme is non-interactive and so the client
%latency is not affected by the server throughput. Compared to the private exact hash
%matching (PEMC) scheme of , our scheme is more than $2\times$ faster with more than two orders of magnitude less communication.  

% a significant burden for storage-limited devices.

Fuzzy matching takes less than 600ms for the fast path check, and 61 seconds for a full check against a blocklist with 1 million items, with false-negative rate of $< 0.08$. The fast path is \textbf{5$\times$} faster than the state-of-the-art PAMC protocol of \citet{kulshrestha2021identifying} with lower overall communications costs by a factor of \textbf{25$\times$}.

%matching scheme takes less than 1.5 seconds to identity a potential match against a million-element blocklist of 256-bit perceptual hash vectors,
% The fine-grained check takes about $127$ seconds \emph{but is only performed if indicated by the coarse-grained check}. The fine-grained check eliminates all false positives without revealing any content to the server. In the average case, our protocol is faster than the private \textit{approximate} hash matching (PAMC) scheme of \citet{kulshrestha2021identifying}, which has a compute time of $27$ seconds and \emph{operates under a more restrictive threat model with a semi-honest server}. 


%In comparison, the private \textit{approximate} hash
%matching (PAMC) scheme of \citet{kulshrestha2021identifying}, which
%has a much higher false negative rate around 17\%, has a compute time of $27$ seconds.  Moreover, our scheme tolerates a more permissive threat model with a malicious server. 


